\documentclass{article}
\usepackage[final]{neurips_2025}
\usepackage[utf8]{inputenc}
\usepackage[T1]{fontenc}
\usepackage{hyperref}
\usepackage{url}
\usepackage{booktabs}
\usepackage{amsfonts}
\usepackage{amsmath}
\usepackage{amssymb}
\usepackage{nicefrac}
\usepackage{microtype}
\usepackage{graphicx}
\usepackage{multirow}
\usepackage{enumitem}

\title{Mitigating Shortcut Reasoning in Language Models: A Gradient-Aware Training Approach}
\author{Author Name \\ Institution \\ \texttt{email@example.com}}
\date{}

\begin{document}
\maketitle

\begin{abstract}
Large language models (LLMs) often rely on reasoning shortcuts---surface pattern matching, answer memorization, and keyword correlations---rather than genuine logical inference. This stems from training signal misalignment: shortcut-encouraging samples efficiently reduce loss but harm generalization under distribution shift. We formulate the problem as \textbf{Shortcut-Aware Reasoning Training (SART)}: identifying training samples that promote shortcut updates and modifying gradient dynamics to emphasize genuine reasoning signals. Our key insight is that shortcut samples produce distinct gradient signatures---low alignment with validation gradients and high concentration on answer tokens. We propose \textbf{ShortcutScore}, a composite metric combining these two signals for sample reweighting, and \textbf{gradient surgery}, which projects harmful gradient directions onto orthogonal subspaces to neutralize shortcut components while preserving reasoning signals. Experiments on mathematical, financial, and causal reasoning benchmarks demonstrate that SART simultaneously improves clean accuracy (+16.2\,pp) and robustness under perturbation (+38.0\,pp) over standard fine-tuning, achieves shortcut detection F1 up to 0.83, and attains the best robustness on multi-feature reasoning tasks, with approximately $2\times$ training overhead. Code: \url{https://github.com/fuyanjie/short-cut-aware-data-centric-reasoning}
\end{abstract}


\input{Introduction2}
\input{Method2}
%=============================================================================
\section{Experiments}
\label{sec:experiments}

We evaluate SART on three synthetic reasoning benchmarks with controlled shortcut injection.

%-----------------------------------------------------------------------------
\subsection{Setup}
\label{sec:exp_setup}

\paragraph{Model and Training.}
We use a GPT-style transformer (6 layers, $d_{\mathrm{model}}=512$, 8 heads, $d_{\mathrm{ff}}=2048$, $\approx$19.0M parameters). All models are trained for 50 epochs with AdamW~\cite{loshchilov2019decoupled}, learning rate $1\times10^{-3}$, weight decay $10^{-4}$, batch size 128, and cosine annealing. For SART: $\alpha=\beta=1.0$, $\tau_A=0.3$, $\tau_R=0.5$, $\lambda=2.0$, $\gamma=0.8$, $\rho=0.7$.

\paragraph{Datasets.}
We construct three synthetic datasets with controlled shortcut injection, each with 10{,}000 train / 2{,}000 val / 3{,}000 test samples. Training labels are 70\% shortcut-rule and 30\% true-rule; validation and test always follow the true rule. The test set is split evenly into \emph{clean} and \emph{perturbed} (shortcut contradicts the correct answer).
\begin{itemize}[leftmargin=*, itemsep=1pt]
    \item \textbf{Math-Arithmetic}: Classify $a+b \geq 10$. Shortcut: $a \geq 5$; true rule: evaluate the full sum.
    \item \textbf{Financial-Analysis}: Determine regulatory compliance from four features. Shortcut: revenue $\geq 5$; true rule: margin $\geq 5$ and debt $< 5$.
    \item \textbf{Causal-Reasoning}: Determine causal direction from observed variables and a confounder. Shortcut: correlation $\geq 5$; true rule: $x \geq 5$ and $z < 3$.
\end{itemize}

\paragraph{Baselines.}
We compare SART against six baselines spanning inference-time, data-centric, loss-based, and distributionally robust approaches:
(1)~\textbf{Standard Fine-Tuning (SFT)}: vanilla cross-entropy, no shortcut mitigation;
(2)~\textbf{Self-Consistency Decoding}~\cite{wang2023selfconsistency}: majority vote over 5 samples at temperature 0.8;
(3)~\textbf{Data Filtering}: remove samples with answer-token confidence $> 0.90$ after a 3-epoch warmup, then retrain;
(4)~\textbf{JTT (Just Train Twice)}~\cite{liu2021just}: identify error-prone samples in a warmup phase and upweight them during retraining;
(5)~\textbf{Focal Loss}~\cite{lin2017focal}: down-weight well-classified examples via $(1-p_t)^\gamma$ weighting;
(6)~\textbf{Group DRO}~\cite{sagawa2020distributionally}: minimize worst-group loss using shortcut/non-shortcut group annotations.

\paragraph{Metrics.}
Clean accuracy, perturbed accuracy (robustness), reasoning consistency (fraction of traces matching ground-truth patterns), and shortcut detection F1 (where applicable).

%-----------------------------------------------------------------------------
\subsection{Main Results}
\label{sec:exp_main}

\begin{table}[t]
\centering
\caption{Main results averaged across three synthetic reasoning datasets. \textbf{Bold}: best per column; \underline{underline}: second best.}
\label{tab:main_results}
\small
\begin{tabular}{lcccc}
\toprule
\textbf{Method} & \textbf{Accuracy} ($\uparrow$) & \textbf{Robustness} ($\uparrow$) & \textbf{Reasoning} ($\uparrow$) & \textbf{SC Det.\ F1} ($\uparrow$) \\
\midrule
Standard Fine-Tuning        & 59.6 & 1.9  & 43.8 & -- \\
Self-Consistency Decoding   & 64.1 & 12.4 & 43.8 & -- \\
Data Filtering              & 68.6 & 20.2 & 57.7 & \textbf{0.82} \\
JTT (Just Train Twice)      & 59.0 & 1.0  & 43.1 & -- \\
Focal Loss                  & 60.2 & 3.8  & 44.8 & -- \\
Group DRO                   & \textbf{78.2} & \textbf{48.2} & \textbf{69.9} & -- \\
\midrule
SART (Ours, Full)           & \underline{75.8} & \underline{39.9} & \underline{66.0} & \underline{0.66} \\
\bottomrule
\end{tabular}
\end{table}

\begin{table}[t]
\centering
\caption{Per-dataset clean and perturbed accuracy (\%). \textbf{Bold}: best; \underline{underline}: second best.}
\label{tab:per_dataset}
\small
\begin{tabular}{llccc}
\toprule
\textbf{Dataset} & \textbf{Method} & \textbf{Acc.\ Clean} ($\uparrow$) & \textbf{Acc.\ Perturbed} ($\uparrow$) & \textbf{Reasoning} ($\uparrow$) \\
\midrule
\multirow{7}{*}{\textsc{Math}}
    & Standard FT          & 75.9 & 0.0  & 75.9 \\
    & Self-Consistency     & 78.3 & 10.7 & 75.9 \\
    & Data Filtering       & 76.9 & 3.0  & 76.9 \\
    & JTT                  & 75.9 & 0.0  & 75.9 \\
    & Focal Loss           & 75.9 & 0.0  & 75.9 \\
    & Group DRO            & \textbf{86.1} & \textbf{51.8} & \textbf{86.1} \\
    & SART (Ours)          & \underline{83.7} & \underline{23.9} & \underline{83.7} \\
\midrule
\multirow{7}{*}{\textsc{Financial}}
    & Standard FT          & 51.5 & 2.9  & 27.8 \\
    & Self-Consistency     & 57.1 & 13.1 & 27.8 \\
    & Data Filtering       & 68.3 & 35.7 & 54.3 \\
    & JTT                  & 50.3 & 1.4  & 26.8 \\
    & Focal Loss           & 53.1 & 8.0  & 31.7 \\
    & Group DRO            & \underline{78.2} & \underline{53.9} & \textbf{64.3} \\
    & SART (Ours)          & \textbf{78.3} & \textbf{58.1} & \underline{56.4} \\
\midrule
\multirow{7}{*}{\textsc{Causal}}
    & Standard FT          & 51.5 & 2.7  & 27.6 \\
    & Self-Consistency     & 56.9 & 13.3 & 27.6 \\
    & Data Filtering       & 60.6 & 21.9 & 41.9 \\
    & JTT                  & 50.8 & 1.7  & 26.5 \\
    & Focal Loss           & 51.5 & 3.3  & 26.8 \\
    & Group DRO            & \textbf{70.3} & \textbf{38.9} & \textbf{59.2} \\
    & SART (Ours)          & \underline{65.5} & \underline{37.8} & \underline{57.9} \\
\bottomrule
\end{tabular}
\end{table}

Table~\ref{tab:main_results} shows that SART achieves 75.8\% clean accuracy (+16.2\,pp over SFT) and 39.9\% robustness (+38.0\,pp), demonstrating that gradient-aware training substantially reduces shortcut reliance. Among all baselines, Group DRO emerges as the strongest competitor (78.2\% accuracy, 48.2\% robustness), leveraging worst-group optimization with known shortcut group annotations.

Several observations are notable.
First, JTT, Focal Loss, and standard fine-tuning cluster at 59--60\% accuracy with near-zero robustness ($\leq$3.8\%), confirming that general-purpose loss modifications alone cannot address systematic shortcut injection.
Second, Data Filtering improves robustness to 20.2\% but at the cost of reduced training data and lower overall accuracy.
Third, Self-Consistency Decoding provides moderate robustness gains (+10.5\,pp) through inference-time diversity, but cannot correct shortcut patterns already encoded during training.

Per-dataset results (Table~\ref{tab:per_dataset}) reveal that \emph{SART achieves the highest accuracy and robustness on Financial-Analysis} (78.3\% / 58.1\%), the most complex multi-dimensional reasoning task, outperforming Group DRO on robustness by 4.2\,pp. This suggests that gradient-based shortcut correction is particularly effective when shortcuts involve complex multi-feature interactions. On Math-Arithmetic and Causal-Reasoning, SART consistently ranks second after Group DRO.

%-----------------------------------------------------------------------------
\subsection{Ablation Study}
\label{sec:exp_ablation}

\begin{table}[t]
\centering
\caption{Ablation of SART components averaged across three datasets. Grad.\ Align.\ measures cosine similarity between training and validation gradients (higher is better).}
\label{tab:ablation}
\small
\begin{tabular}{lccc}
\toprule
\textbf{Configuration} & \textbf{Accuracy} ($\uparrow$) & \textbf{Robustness} ($\uparrow$) & \textbf{Grad.\ Align.} ($\uparrow$) \\
\midrule
SFT (Baseline)              & 59.6 & 1.9   & $-0.07$ \\
\midrule
Reweighting Only            & 64.7 & 13.1  & -- \\
Gradient Surgery Only       & \textbf{81.0} & \textbf{51.3}  & -- \\
\textbf{SART (Full)}        & 75.8 & 39.9  & $\mathbf{+0.10}$ \\
\bottomrule
\end{tabular}
\end{table}

We isolate each component's contribution (Table~\ref{tab:ablation}). Both components individually improve over SFT: Reweighting provides +5.1\,pp accuracy and +11.2\,pp robustness; Gradient Surgery achieves +21.4\,pp accuracy and +49.4\,pp robustness, establishing it as the primary mechanism for shortcut correction. By projecting gradients away from shortcut-correlated directions, Gradient Surgery directly modifies the optimization trajectory, while Reweighting adjusts sample influence in the loss landscape.

The full method achieves the best gradient alignment ($+0.10$ vs.\ $-0.07$ for SFT), indicating the most thorough correction of shortcut-correlated gradient components. Per-dataset analysis reveals complementary behavior: on Math-Arithmetic, the combination achieves synergistic gains ($83.7\%/23.9\%$) over Gradient Surgery alone ($78.0\%/15.3\%$), yielding +5.7\,pp accuracy and +8.6\,pp robustness. On the multi-feature Financial and Causal tasks, Gradient Surgery alone provides particularly strong correction, suggesting that when gradient-level modification is sufficiently aggressive, additional sample-level reweighting may introduce interference. This points to future work on adaptive combination strategies.

%-----------------------------------------------------------------------------
\subsection{ShortcutScore Validation}
\label{sec:exp_detection}

\begin{figure}[t]
    \centering
    \includegraphics[width=\textwidth]{results/figure3.png}
    \caption{Empirical validation of \textsc{ShortcutScore}. (a) Score vs.\ true shortcut rate (Pearson $r=0.67$). (b) Gradient alignment distribution: shortcut samples (red) skew negative; non-shortcut (green) cluster near zero. (c) Exponential reweighting with $\lambda=2.0$.}
    \label{fig:analysis}
\end{figure}

\textsc{ShortcutScore} achieves average shortcut detection F1 of 0.66 across three datasets, with per-task F1 ranging from 0.56 (Financial) to 0.83 (Math). On Math-Arithmetic, \textsc{ShortcutScore} (F1$=$0.83) substantially outperforms confidence-based filtering (F1$=$0.70), confirming that gradient structure captures shortcut reliance that output confidence cannot detect---shortcut samples are correctly labeled and confidently predicted, making them indistinguishable from genuinely easy samples by probability alone. On multi-feature tasks (Financial, Causal), confidence-based filtering achieves higher detection F1, suggesting that the simpler threshold-based heuristic may be more effective when shortcuts involve correlated feature groups. Figure~\ref{fig:analysis}(a) confirms a positive monotonic correlation between $S(s)$ and true shortcut rate ($r=0.67$). The gradient alignment distribution (Figure~\ref{fig:analysis}b) shows clear separation: shortcut samples skew toward negative alignment, while non-shortcut samples cluster near zero.

\subsection{Computational Cost}
\label{sec:exp_efficiency}

The full pipeline completes in under 3 minutes per dataset on a single NVIDIA H100 NVL GPU. SART incurs approximately $2\times$ training overhead versus SFT, attributable to per-sample gradient computation for \textsc{ShortcutScore} and gradient projection during surgery. Validation gradients are recomputed periodically rather than per-step, keeping overhead manageable. This cost can be further reduced via gradient checkpointing or random projections for larger models.



\section{Related Work}

\textbf{Language model training and reasoning.}
Current training paradigms predominantly optimize for answer correctness, inadvertently encouraging shortcut exploitation~\cite{Author_Year_KeyWord}. Chain-of-Thought training~\cite{wei2022chain,kojima2022large} encourages intermediate reasoning traces but does not modify training dynamics, allowing shortcuts to persist in learned parameters---models can generate plausible reasoning steps that still rely on shortcuts. Self-Consistency Decoding~\cite{wang2022self} samples multiple reasoning paths to detect inconsistencies at inference time but cannot correct shortcut patterns already encoded during training, leaving the core training process unchanged. Data Filtering and Curriculum Learning~\cite{bengio2009curriculum} remove noisy or difficult samples but miss correctly labeled shortcut samples that appear indistinguishable from high-quality data. RLHF~\cite{ouyang2022training} optimizes reward signals based on answer correctness, not reasoning validity, allowing shortcuts to persist as long as they produce correct answers on the training distribution. These approaches demonstrate the inability to jointly address shortcut detection from gradient behavior and dynamic suppression during training.
%
\textbf{Gradient-based optimization for robustness.}
Standard gradient methods cannot distinguish reasoning-improving gradients from shortcut-reinforcing ones~\cite{Author_Year_KeyWord}. Prior gradient modification techniques~\cite{Author_Year_KeyWord,yu2020gradient} improve general robustness but lack principled mechanisms to identify \emph{which} gradients are detrimental to reasoning and \emph{why}---they modify gradients uniformly without targeting shortcut-specific patterns. Training signal adjustment methods~\cite{Author_Year_KeyWord} do not incorporate gradient-based detection of shortcut samples and cannot quantify the ``shortcut-ness'' of a sample based on its gradient behavior. Our work uniquely combines gradient alignment and answer-gradient concentration into a unified ShortcutScore, enabling targeted detection and suppression of shortcut-promoting training signals---a more principled approach than uniform gradient modification or sample removal.


\section{Conclusion}

We presented SART, a gradient-aware training framework that addresses shortcut reasoning in LLMs through two complementary mechanisms: ShortcutScore-based sample reweighting and gradient surgery. Our key finding is that shortcut samples produce distinct gradient signatures---low validation alignment and high answer-token concentration---enabling principled detection (F1 up to 0.83) that captures shortcut reliance invisible to confidence-based filtering. Experimental results demonstrate simultaneous improvements in accuracy (+16.2\,pp) and robustness (+38.0\,pp) over standard fine-tuning, with SART achieving the strongest robustness on the most complex multi-feature reasoning task. Ablation analysis reveals that gradient surgery is the primary mechanism, with the full method achieving the best gradient alignment (+0.10 vs.\ $-$0.07 for SFT). The framework operates purely on training dynamics, requires no architectural modifications, and incurs approximately $2\times$ overhead. Theoretically, this work advances the understanding of gradient alignment and reasoning signal correction, offering new insights into LLM training dynamics. Practically, the implications extend to domains such as finance and scientific discovery, where enhanced model robustness leads to more reliable decision-making. Future work will explore adaptive combination strategies and extend SART to larger-scale models and diverse reasoning domains.


\bibliographystyle{plain}
\bibliography{references}

\end{document}
